\section{Aproximaciones}

jeje

\subsection{Teoría}
. trade-off entre tiempo polinomial y optimalidad garantizada (elk r(A))


\subsection{Análisis del Algoritmo del Gringo Hinz}
. Demostración de validez, complejidad O(N.M)
. demostracion matemática de la cota r(A)


\subsubsection{Propuesta de Aproximación Propia}
. el por qué greedy + diseño determinístico
. pseudocodigo
. justificacion de la complejidad


\subsection{Pensamiento Algorítmico en la Vida Real}
. Conexión entre los 4 temas (Herbert Simon y la Racionalidad Limitada)
    . saber cuándo podar un camino con Backtracking y cuándo preferir una aproximación a tiempo frente a la parálisis por análisis).

    
tengo pensado:
En la vida cotidiana y en la ingeniería nos topamos con problemas "NP-Completos" (decisiones con miles de variables interconectadas donde buscar la perfección absoluta toma tiempo infinito). 
Ahí aplicamos los 4 pilares: Reducciones (reconocer que un problema nuevo es en realidad uno viejo disfrazado), Backtracking (explorar caminos pero tener la madurez de podar y retroceder apenas vemos que una rama no supera lo que ya logramos), Programación Lineal (optimizar nuestros objetivos dentro de restricciones rígidas), y Aproximaciones (el principio de Racionalidad Limitada o Satisficing del Premio Nobel y Premio Turing Herbert Simon: cuando encontrar el óptimo exacto es computacionalmente intratable, una buena aproximación a tiempo es infinitamente mejor que la perfección que llega tarde).